% ============================================================
% IEEE TALE 2026 — Conference Paper
% Generative AI in Assessment: A Systematization of Knowledge
% Template: IEEEtran (conference mode)
% ============================================================
\documentclass[conference]{IEEEtran}

% ── Packages ──────────────────────────────────────────────────────────────────
\usepackage{cite}           % Sort and compress citation numbers
\usepackage{amsmath}
\usepackage{booktabs}       % \toprule, \midrule, \bottomrule
\usepackage{array}          % Column type control
\usepackage{url}
\usepackage[hidelinks]{hyperref}
\usepackage{graphicx}
\usepackage{xcolor}
\usepackage{multirow}
\usepackage{tabularx}       % Auto-width columns
\usepackage{ragged2e}       % \RaggedRight in table cells

% Fix IEEEtran + hyperref interaction
\makeatletter
\AtBeginDocument{%
  \let\IEEEproof\proof
  \let\IEEEendproof\endproof
}
\makeatother

% Slightly tighter spacing in tables
\renewcommand{\arraystretch}{1.15}

% ── Metadata ──────────────────────────────────────────────────────────────────
\title{Generative AI in Assessment within Engineering, STEM, and Computing Education:
       A Systematization of Knowledge}

\author{%
  \IEEEauthorblockN{Sian Lerk Lau}
  \IEEEauthorblockA{%
    \textit{School of Mathematical and Computer Sciences}\\
    Sunway University\\
    Petaling Jaya, Malaysia\\
    sianlunl@sunway.edu.my
  }
}

% ── Begin document ─────────────────────────────────────────────────────────────
\begin{document}

\maketitle

% ── Abstract ──────────────────────────────────────────────────────────────────
\begin{abstract}
The public release of ChatGPT in November 2022 catalysed a rapid expansion of scholarship on
generative AI~(GenAI) in educational assessment, yet the literature has outpaced synthesis
efforts. This paper presents a Systematization of Knowledge~(SoK) review drawing on 641
peer-reviewed papers sourced from five databases (Scopus, Web of Science, IEEE Xplore, ACM
Digital Library, and ERIC), spanning 2022 to mid-2026, with bibliometric analysis covering
2,122 deduplicated records. A dual-pipeline methodology was employed: bibliometric network
analysis and keyword-based thematic coding of 617 full-text PDFs. Annual output grew from 114
papers in 2023 to 1,054 in 2025 (CAGR~296.6\%), with publication dominated by ASEE, ACM
SIGCSE, and FIE proceedings. Keyword co-occurrence analysis identifies four primary clusters:
Disciplinary Contexts ($n{=}1{,}681$ co-occurrences), GenAI/Assessment Discourse
($n{=}1{,}662$), LLM Applications in Education ($n{=}1{,}417$), and CS/Computing Pedagogy
($n{=}987$). Full-text taxonomy coding across seven categories reveals that ethics and policy
concerns~(90.3\%), perception and affect~(92.7\%), and review methodology~(93.5\%) dominate,
while automated grading~(47.0\%) and authentic assessment redesign~(44.9\%) remain
comparatively thin. Evidence quality mapping reveals a consistent structural gap: despite high
publication volume, fewer than one in ten papers reports a controlled evaluation of an
implemented intervention. These findings reveal a field characterised by diagnostic breadth
but insufficient operationalisation, and identify automated grading validity, AI-integrated
assessment design, and cross-national equity studies as highest-priority research directions.
\end{abstract}

\begin{IEEEkeywords}
generative AI; assessment; engineering education; STEM; computing education; systematization
of knowledge; large language models; evidence synthesis
\end{IEEEkeywords}

% ── I. Introduction ───────────────────────────────────────────────────────────
\section{Introduction}

The public release of ChatGPT in November 2022 marked a qualitative shift in the relationship
between artificial intelligence and education. Unlike preceding AI tools that augmented
specific narrow tasks, large language models~(LLMs) capable of generating coherent academic
prose, executable code, and plausible solutions to engineering problems challenged foundational
assumptions about what student-produced work could and should demonstrate. Within months,
academic institutions worldwide were grappling with questions that traditional assessment
frameworks had not been designed to answer: How do we verify authorship? What constitutes
authentic demonstration of competence? How should formative feedback be redesigned when
students can generate instant, detailed responses from an AI?

These questions are particularly acute in engineering, STEM, and computing disciplines, where
the outputs most readily produced by GenAI tools---code, calculations, and structured technical
reports---overlap substantially with standard assessment artefacts. The scholarly response has
been rapid: annual publication output grew from 114 papers in 2023 to 1,054 in 2025, a rate
that outpaces even the GenAI adoption curve itself. Without systematic mapping, the field risks
duplicating effort, pursuing fashionable sub-themes at the expense of critical gaps, and failing
to translate a growing evidence base into actionable guidance for practitioners.

This paper addresses that need with a Systematization of Knowledge~(SoK) review---a methodology
that produces structured, reproducible overviews of rapidly developing fields by mapping not only
\textit{what} is known, but \textit{how well} it is known and \textit{what remains unknown}.
Unlike a conventional systematic review that synthesises around a pre-specified research question,
a SoK explicitly characterises the shape of knowledge: its conceptual clusters, thematic emphases,
evidence quality distribution, and structural gaps. Our primary contribution is a structured account
of what nearly four years of rapid scholarship has established, with what confidence, and where the
field must go next. The corpus and pipeline code are deposited in an open repository~\cite{sok_repo}.

Four research questions guide this SoK. \textbf{RQ1}: What is the bibliometric landscape of
GenAI-in-assessment research within engineering, STEM, and computing education from 2022 to
mid-2026? \textbf{RQ2}: What thematic clusters emerge from keyword co-occurrence analysis?
\textbf{RQ3}: What is the distribution and evidence quality of research across a seven-category
taxonomy? \textbf{RQ4}: What cross-cutting patterns and priority gaps emerge from integrating
both pipelines? Section~\ref{sec:method} describes the methodology; Section~\ref{sec:corpus}
presents corpus characterisation; Section~\ref{sec:taxonomy} delivers the SoK taxonomy synthesis;
Sections~\ref{sec:crosscut} and~\ref{sec:gaps} present cross-cutting findings and the gap map;
Sections~\ref{sec:implications}--\ref{sec:conclusion} cover implications, limitations, and
conclusions.

% ── II. Methodology ───────────────────────────────────────────────────────────
\section{Methodology}
\label{sec:method}

\subsection{Search Strategy}

A multi-database search was conducted in June 2026 across Scopus, Web of Science~(WoS), IEEE
Xplore, ACM Digital Library, and ERIC. The search string combined generative AI terms
(\textit{generative AI, ChatGPT, large language models, LLM, GPT}) with assessment-related
terms (\textit{assessment, grading, evaluation, academic integrity, feedback}) and discipline
scope (\textit{engineering education, STEM, computing education, computer science education}).
The temporal window was 2022 to mid-2026. No language restriction was applied at search; all
non-English papers were excluded at full-text review.

\subsection{PRISMA-Aligned Screening}

A total of 2,766 records were retrieved prior to deduplication. After removing 604
cross-database duplicates (DOI- and title-based matching), 2,162 unique records were screened
on title and abstract. Inclusion required substantive engagement with GenAI in at least one of:
(i)~any component of the assessment cycle; (ii)~assessment redesign in response to GenAI;
(iii)~academic integrity implications; or (iv)~student or instructor perceptions of GenAI in
assessment contexts. After two-stage screening (automated keyword screening followed by manual
review of 59 uncertain cases), 641 papers were retained:
\[
2{,}766 \xrightarrow{-604} 2{,}162 \xrightarrow{-1{,}517} 645 \xrightarrow{-4} 641
\]

\subsection{Dual-Pipeline Analysis}

\textit{Pipeline~1} (bibliometric analysis) used the \texttt{bibliometrix} R package~\cite{aria2017}
on a merged five-database dataset. IEEE Xplore, ACM, and ERIC BibTeX exports were first normalised
to Scopus-compatible format via a custom Python script before loading. Author affiliation data absent
from IEEE/ACM/ERIC exports was retrospectively enriched via the OpenAlex API, achieving 96.9\%
coverage (947/977 records with DOIs). Analyses included annual production trends, venue and author
productivity, country collaboration networks, and keyword co-occurrence mapping.

\textit{Pipeline~2} (content analysis) extracted full text from 617 PDFs using PyMuPDF, then
applied keyword-based thematic coding across seven predefined categories. Papers could be assigned
to multiple categories, yielding co-occurrence rates rather than exclusive classification. The two
pipelines are analytically independent: convergence between them strengthens findings; divergence
identifies areas requiring deeper investigation.

% ── III. Corpus Characterisation ──────────────────────────────────────────────
\section{Corpus Characterisation (RQ1 and RQ2)}
\label{sec:corpus}

\subsection{Growth and Scope}

The five-database corpus of 2,122 records reveals a dramatic growth trajectory (Table~\ref{tab:annual}).
Annual output rose from a near-zero baseline of 2 papers in 2022 to 1,054 in 2025, a compounded
annual growth rate of 296.6\%. The 2026 partial count of 495 papers to mid-year suggests the
trajectory is sustained. The document average age of 1.09 years confirms the extreme recency of the
field, and the average citation count of 2.38 per document reflects the characteristic low-citation
state of an emerging literature yet to accumulate referencing density. Conference papers dominate
(71.3\%; $n{=}1{,}513$), with articles (17.0\%) and proceedings papers (9.0\%) comprising the
remainder.

\begin{table}[ht]
\caption{Annual Publication Volume --- All Five Databases ($N = 2{,}122$)}
\label{tab:annual}
\centering
\begin{tabular}{crrr}
\toprule
\textbf{Year} & \textbf{Papers} & \textbf{YoY Growth} & \textbf{Cumulative} \\
\midrule
2022 & 2     & ---       & 2     \\
2023 & 114   & +5,600\%* & 116   \\
2024 & 457   & +301\%    & 573   \\
2025 & 1,054 & +131\%    & 1,627 \\
2026\dag & 495 & (partial) & 2,122 \\
\bottomrule
\multicolumn{4}{l}{\footnotesize * 2022 baseline of 2 inflates figure. CAGR (2022--2025)~=~296.6\%.}\\
\multicolumn{4}{l}{\footnotesize \dag To mid-2026 only.}
\end{tabular}
\end{table}

\subsection{Venue and Author Landscape}

The corpus spans 777 unique sources, yet publication is concentrated: the top five venues account
for 11.3\% of records. ASEE Annual Conference proceedings lead with 61 papers, followed by the
ACM SIGCSE Technical Symposium (47 and 45 papers across its 2025 and 2026 editions),
Frontiers in Education Conference~(43), and Lecture Notes in Networks and Systems~(43). The
ACM SIGCSE symposium---substantially undercounted in Scopus/WoS-only analyses---represents a
key computing education venue that the five-database approach recovers. The most productive
author is Paul Denny~(51 papers), with Juho Leinonen~(35), Stephen MacNeil~(24), James
Prather~(21), and Chen~Y~(20) following. The most-cited paper is Bahroun et al.~\cite{bahroun2023}
in \textit{Sustainability} (a broad ChatGPT-in-education review; 361 citations), with Nikolic
et al.~\cite{nikolic2023} in the \textit{European Journal of Engineering Education}~(173
citations) as the most-cited engineering-specific work.

\subsection{Geographic Distribution}

The United States accounts for the largest share of publications~(18.8\%), followed by
China~(15.2\%), India~(5.0\%), the United Kingdom~(4.0\%), and Australia~(3.9\%). International
co-authorship stands at 9.9\% of documents---substantially higher than earlier Scopus/WoS-only
estimates once IEEE and ACM affiliation data is enriched via OpenAlex. The mean co-author count
of 3.89 per document is consistent with engineering education norms. Citation impact by country
shows UAE~(avg.\ 81.0 citations/paper), Sweden~(21.8), and New Zealand~(16.9) producing the
most highly cited papers, suggesting foundational contributions from smaller research communities.

\subsection{Keyword Co-occurrence Clusters (RQ2)}

Keyword co-occurrence mapping of the merged five-database corpus identified four primary clusters
(Table~\ref{tab:clusters}). Total co-occurrences sum across author keyword~(DE) and Keywords
Plus~(ID) fields within each cluster.

\begin{table}[ht]
\caption{Keyword Co-occurrence Clusters}
\label{tab:clusters}
\centering
\begin{tabular}{clr}
\toprule
\textbf{\#} & \textbf{Cluster Label} & \textbf{Co-occ.} \\
\midrule
1 & Disciplinary Contexts       & 1,681 \\
2 & GenAI / Assessment Discourse & 1,662 \\
3 & LLM Applications in Education & 1,417 \\
4 & CS / Computing Pedagogy     & 987   \\
\bottomrule
\end{tabular}
\end{table}

Cluster~1 (\textit{Disciplinary Contexts}; 1,681) functions as a structural backdrop, linking
disciplinary homes (engineering, CS, STEM), pedagogical structures (e-learning, educational
technology), and normative frameworks (ethics) to the GenAI phenomenon. Cluster~2
(\textit{GenAI/Assessment Discourse}; 1,662) is the dominant discursive hub---combining the
technology labels (\textit{generative ai, chatgpt}) with assessment concerns (\textit{academic
integrity, higher education})---representing the problem space of the literature. Cluster~3
(\textit{LLM Applications}; 1,417) groups LLM capabilities with educational applications: prompt
engineering, automated grading, feedback generation, and personalised learning. Cluster~4
(\textit{CS/Computing Pedagogy}; 987) reflects the sub-field most active in LLM research---
introductory computing education, CS1 course pedagogy, and programming assessment---whose
distinctness from Cluster~3 reflects a disciplinary framing rather than a general educational
technology one.

% ── IV. Taxonomy and Evidence Synthesis ───────────────────────────────────────
\section{Taxonomy and Evidence Synthesis (RQ3)}
\label{sec:taxonomy}

Full-text extraction and thematic coding of 617~PDFs produced a seven-category taxonomy
(Table~\ref{tab:taxonomy}). For each category we summarise \textit{what the literature
establishes}, \textit{the quality of evidence}, and \textit{the critical gap}. Categories are
ordered from least (A) to most (G) prevalent to highlight where thinner coverage coincides with
higher practical urgency.

\begin{table*}[ht]
\caption{Seven-Category SoK Taxonomy --- Distribution and Evidence Quality ($N = 617$ coded papers)}
\label{tab:taxonomy}
\centering
\begin{tabular}{clrrp{6.8cm}}
\toprule
\textbf{Cat.} & \textbf{Theme} & \textbf{Papers} & \textbf{\%} & \textbf{Evidence Quality} \\
\midrule
A & Automated Grading              & 290 & 47.0\% & Moderate --- small-$n$ empirical studies; limited validity analysis \\
B & Authentic Assessment Redesign  & 277 & 44.9\% & Low --- predominantly conceptual/framework; few evaluations \\
C & Academic Integrity             & 430 & 69.7\% & Moderate --- detection studies; cross-sectional policy surveys \\
D & Formative and Adaptive Assessment & 439 & 71.2\% & Moderate--High --- controlled studies in ITS and CS1 contexts \\
E & Ethics and Policy              & 557 & 90.3\% & Low --- framework/opinion dominant; sparse empirical work \\
F & Perception and Affect          & 572 & 92.7\% & Moderate --- survey-heavy; largely cross-sectional \\
G & Review / Methodology           & 577 & 93.5\% & N/A --- methodological marker; near-universal rate \\
\bottomrule
\multicolumn{5}{l}{\footnotesize Categories are not mutually exclusive; percentages reflect papers assigned to each.}
\end{tabular}
\end{table*}

\subsection{Category A: Automated Grading (47.0\%; $n = 290$)}

\textit{What the literature establishes.}
LLMs can evaluate programming assignments, essays, and short structured answers with moderate
reliability. Studies using GPT-4 class models report human--AI grader agreement rates of
approximately 70--85\% for rubric-based, low-ambiguity tasks and substantially lower agreement
for open-ended design reasoning. AI feedback tends toward verbosity and generality unless prompts
are calibrated with rubric-specific criteria. Concerns about hallucinated justifications in written
feedback remain active. The most prominent research cluster in this area (Denny, Leinonen, MacNeil,
Savelka, Prather) focuses on code grading and code explanation tasks~\cite{macneil2023,savelka2023}.

\textit{Evidence quality: Moderate.}
The empirical base includes human--AI agreement studies against instructor grades as ground truth,
but most studies use convenience samples from single courses ($n < 150$), deploy controlled scenarios
rather than live submissions, and measure agreement without addressing validity or fairness.

\textit{Critical gap.}
No large-scale validity or fairness study has been published for automated grading in engineering
or computing education. The field lacks agreed thresholds for acceptable human--AI agreement for
formative versus summative purposes. Longitudinal data on student learning outcomes when AI grading
replaces human grading is entirely absent.

\subsection{Category B: Authentic Assessment Redesign (44.9\%; $n = 277$)}

\textit{What the literature establishes.}
Two broad design philosophies have emerged. \textit{AI-resistant design} employs formats that GenAI
cannot plausibly complete---synchronous in-person examinations, viva voce oral defences, and
reflective journals requiring verifiable personal experience. \textit{AI-integrated design} explicitly
incorporates AI use, assessed on process rather than product. Artefact-plus-reflection formats,
process portfolios, and live demonstrations are the most commonly proposed implementations. Bloom's
taxonomy reframing is a recurring motif: shift assessment emphasis toward synthesis, evaluation, and
creative design that AI performs less well than recall and application.

\textit{Evidence quality: Low.}
The dominant publication type is conceptual or framework-based: proposed redesign models, opinion
pieces, and institutional case narratives. Controlled comparisons between AI-resistant and AI-integrated
designs---and between either and baseline assessment---are almost entirely absent.

\textit{Critical gap.}
No published study directly compares learning outcomes under AI-resistant versus AI-integrated
assessment in engineering or computing education. Scalability of high-authenticity formats in large
cohorts remains empirically unaddressed.

\subsection{Category C: Academic Integrity (69.7\%; $n = 430$)}

\textit{What the literature establishes.}
AI detection tools (Turnitin AI, GPTZero, Copyleaks) show false positive rates of 8--20\% in
controlled studies, with substantially higher misclassification for non-native English speakers.
Detection-as-policy is widely critiqued as an arms race. Drawing on contract cheating research
predating GenAI, the more productive literature finds that institutional culture, task design, and
policy communication matter more than technical detection. Student attitudes toward disclosure are
more cooperative than commonly assumed when clear guidance and proportionate penalties exist.

\textit{Evidence quality: Moderate, but uneven.}
Detection tool studies have quantitative metrics but suffer from controlled-condition design (ground
truth AI text versus human text, rather than real hybrid-use submissions). Policy surveys provide
cross-institutional coverage but are cross-sectional. Student self-report data on AI use is susceptible
to social desirability bias.

\textit{Critical gap.}
No longitudinal dataset tracks actual integrity violation rates before and after institutional policy
interventions. The differential impact of AI detection systems on non-native English speakers has
been identified but not addressed with intervention research.

\subsection{Category D: Formative and Adaptive Assessment (71.2\%; $n = 439$)}

\textit{What the literature establishes.}
This is the category with the most developed empirical base, drawing on decades of intelligent
tutoring system~(ITS) research. LLM integration into formative loops enables personalised hints,
Socratic questioning, and immediate explanatory feedback on code, mathematics, and engineering
problem steps. CS1 and introductory programming contexts dominate. Controlled studies report
measurable short-term learning gains on targeted skills. Concurrently, the literature identifies a
risk of learned helplessness when students receive immediate AI hints without productive struggle.

\textit{Evidence quality: Moderate to High.}
The ITS lineage provides a rigorous methodological baseline; several papers report controlled
experiments with pre-post designs or comparable control groups. The limitation is ecological validity:
most studies are in introductory programming with motivated participants, and study durations rarely
exceed one semester.

\textit{Critical gap.}
Long-term retention and transfer from LLM-based formative feedback are unstudied. Comparative
evaluation of human-in-the-loop versus fully automated formative systems is needed. Applications
beyond introductory computing---into engineering design critique and upper-division STEM---represent
an underexplored deployment space.

\subsection{Category E: Ethics and Policy (90.3\%; $n = 557$)}

\textit{What the literature establishes.}
Near-universal engagement with ethics and policy is a defining characteristic of this field. Recurring
concerns include bias and fairness in AI grading, transparency, data privacy, environmental costs of
LLM deployment, and equity of access to premium AI tools. Multiple frameworks have been proposed at
institutional, national, and international levels (UNESCO, ACM, IEEE). The equity dimension is
particularly prominent: unequal access to GPT-4 versus free-tier models creates assessment conditions
that differ across socioeconomic groups.

\textit{Evidence quality: Low.}
The vast majority of ethics and policy papers are conceptual, framework-based, or opinion pieces.
Empirical ethics research---measuring the actual magnitude of AI grading bias or quantifying the
equity impact of tool access differentials---is sparse.

\textit{Critical gap.}
Policy effectiveness research is almost entirely absent: no published study tracks whether a specific
institutional AI policy (ban, allow with disclosure, AI-integrated) produces different learning or
integrity outcomes over time. Cross-national comparative policy analysis has not appeared in this
corpus despite frequent advocacy for it.

\subsection{Category F: Perception and Affect (92.7\%; $n = 572$)}

\textit{What the literature establishes.}
Students are, on average, more positive toward GenAI in assessment contexts than instructors, but
attitudes are neither uniform nor stable. A consistent finding is attitudinal bifurcation: enthusiasm
for efficiency gains coexists with anxiety about dependency and skills atrophy. Cultural variation is
significant---East Asian student samples show stronger concern about integrity norms; Western samples
more readily frame GenAI as a neutral tool. AI-related anxiety correlates with lower academic
self-efficacy. Initial enthusiasm moderates with sustained use as limitations become apparent.

\textit{Evidence quality: Moderate.}
Approximately 76\% of perception papers use questionnaire methods. The literature is overwhelmingly
cross-sectional. Validated instruments are emerging (AI Attitude Scale variants) but no standardised
instrument has achieved consensus, making cross-study comparison unreliable.

\textit{Critical gap.}
Longitudinal tracking of attitude change across a full academic year---as students and instructors
gain sustained experience---is absent. The relationship between perception and learning outcomes is
essentially unstudied. Cross-cultural comparative studies using equivalent instruments are needed.

\subsection{Category G: Review / Methodology (93.5\%; $n = 577$)}

This category warrants special interpretation. Its near-universal prevalence reflects the breadth of
keyword matching rather than a distinct substantive research area. Substantively, the most important
observation is that the field is self-aware of its proliferation---multiple bibliometric and systematic
reviews have been published within the corpus window itself. This recursive dynamic points to a
community actively constructing knowledge-management infrastructure in real time. No dominant study
design has emerged; the corpus mixes surveys, case studies, quasi-experiments, and technology
demonstrations, with randomised controlled trials representing fewer than 5\% of empirical papers.

% ── V. Cross-Cutting Findings ─────────────────────────────────────────────────
\section{Cross-Cutting Findings (RQ4)}
\label{sec:crosscut}

\textbf{The diagnosis-to-intervention gap.}
The most consistent finding across both pipelines is a structural gap between diagnostic and
interventional research. Categories with highest coverage (E, F) are precisely those characterised
by the lowest evidence quality. Categories with highest practical urgency and most developed empirical
traditions (A, D) show lower overall prevalence but stronger study designs. This inversion---abundant
low-quality coverage of concerns, sparse high-quality evaluation of solutions---is characteristic of
fields responding to rapid disruptive change but creates a genuine gap between the literature's volume
and its actionability.

\textbf{Methodological homogeneity.}
Across the taxonomy, survey-based and conceptual/framework methodologies dominate. This produces a
literature rich in documented perceptions and proposed frameworks but thin on controlled evaluation
of implementation outcomes. Computing education shows the most methodological diversity and the most
controlled studies, reflecting its longer tradition of education research. The field would benefit from
longer-term naturalistic studies, quasi-experimental designs exploiting natural variation in institutional
policies, and mixed-methods longitudinal work linking attitude data to outcome data.

\textbf{Disciplinary unevenness.}
Computing education drives a disproportionate share of the empirical and intervention literature,
reflecting ACM Digital Library coverage, methodological tradition, and the direct alignment between LLM
capabilities and computing assessment artefacts. Engineering and broader STEM contexts are
underrepresented in the high-impact citation tier despite inclusion in the search strategy. The Asia-Pacific
region, though contributing significantly to publication volume, is underrepresented in the most-cited work.

\textbf{The AI-proofing versus AI-integration debate.}
A recurring conceptual tension---particularly pronounced in categories~B and~E---is whether educational
response to GenAI should focus on AI-resistant assessment or redesigning assessments to incorporate AI
as a legitimate tool. The literature does not resolve this tension empirically; both positions are
represented by framework papers and institutional case studies without comparative outcome data.

% ── VI. Gap Map ───────────────────────────────────────────────────────────────
\section{Gap Map and Research Agenda}
\label{sec:gaps}

Table~\ref{tab:gaps} maps priority research gaps identified through the SoK synthesis. Gaps are
categorised by urgency (the practical or theoretical cost of the knowledge absent) and feasibility
(whether the gap can be addressed with currently available methodology and data). High
urgency / high feasibility gaps represent the most actionable priorities for the next research cycle.

\begin{table*}[ht]
\caption{Priority Research Gap Map}
\label{tab:gaps}
\centering
\begin{tabular}{p{2.4cm}p{8.2cm}cc}
\toprule
\textbf{Category} & \textbf{Gap} & \textbf{Urgency} & \textbf{Feasibility} \\
\midrule
Automated Grading (A)
  & Validity and fairness standards for AI grading across demographic groups
  & High & High \\
Automated Grading (A)
  & Longitudinal learning outcomes under AI vs.\ human grading
  & High & Medium \\
Authentic Redesign (B)
  & Controlled comparison of AI-resistant vs.\ AI-integrated designs on learning outcomes
  & High & Medium \\
Authentic Redesign (B)
  & Scalability of oral examination formats in large engineering cohorts
  & High & High \\
Academic Integrity (C)
  & Longitudinal integrity violation rates before/after institutional policy changes
  & High & Medium \\
Academic Integrity (C)
  & Fairness of AI detection systems across student language backgrounds
  & High & High \\
Formative / Adaptive (D)
  & Long-term retention and transfer from LLM-based formative feedback
  & High & Medium \\
Formative / Adaptive (D)
  & LLM formative assessment in engineering design and upper-division STEM
  & Medium & High \\
Ethics / Policy (E)
  & Policy effectiveness evaluation: which institutional AI policies improve outcomes?
  & High & Medium \\
Ethics / Policy (E)
  & Cross-national comparative analysis of institutional policy responses
  & Medium & High \\
Perception / Affect (F)
  & Longitudinal attitude change over a full academic year of GenAI use
  & Medium & High \\
Perception / Affect (F)
  & Cross-cultural comparative studies with standardised instruments
  & Medium & High \\
Cross-cutting
  & Equity impact of unequal AI tool access on assessment outcomes
  & High & Medium \\
Cross-cutting
  & Non-Western educational context representation in empirical studies
  & Medium & High \\
\bottomrule
\end{tabular}
\end{table*}

% ── VII. Implications ─────────────────────────────────────────────────────────
\section{Implications}
\label{sec:implications}

\subsection{For Assessment Designers and Instructors}

The taxonomy synthesis supports four evidence-grounded design principles.
\textit{First}, formative AI feedback~(Category~D) has the strongest empirical base and can be
adopted with confidence for introductory-level tasks in computing and STEM, provided that
productive struggle time is preserved and learned helplessness is monitored.
\textit{Second}, process-based assessment components---reflection prompts, design process
documentation, viva elements---are the most consistently advocated AI-integration strategy~(B),
though instructors adopting these designs should document outcomes to contribute to the emerging
evidence base.
\textit{Third}, AI detection tools should not be the primary integrity mechanism~(C); task
redesign, policy clarity, and disclosure culture are more effective.
\textit{Fourth}, student perception data~(F) confirms that clear, proportionate, and
student-co-designed AI use policies reduce anxiety and improve disclosure rates.

\subsection{For Institutions and Policymakers}

The near-universal engagement with ethics and policy~(E) without accompanying effectiveness data
represents an accountability gap: institutions are publishing AI policies without evaluating their
outcomes. This SoK recommends that institutions treat their AI policies as researchable
interventions---documenting the policy, rationale, and measured outcomes over at least one academic
year. The geographic concentration of the literature means that policies developed in large US
research university contexts may not transfer to smaller institutions, community colleges, or
non-Western systems. IEEE TALE, with its Asia-Pacific focus, is well positioned to address this
gap. The equity concern about unequal access to premium AI tools requires immediate institutional
response: where AI tool access is permitted in assessment, institutions should either provide
equitable access or explicitly account for tool access differences in assessment design.

\subsection{For Researchers}

Table~\ref{tab:gaps} identifies fourteen priority research directions. The highest combined urgency
and feasibility targets are: automated grading validity and fairness studies~(A), detection system
fairness for non-native English speakers~(C), scalability of oral examination formats in large
cohorts~(B), and longitudinal attitude studies~(F). These can all be addressed with existing
methodology and data sources. The field would benefit from pre-registration of studies addressing
high-stakes outcomes (integrity violation rates, learning gains from formative feedback). Researchers
in non-Western contexts---particularly Asia-Pacific, South Asia, and Latin America---are encouraged
to contribute to the comparative literature that is currently underrepresented in this corpus.

% ── VIII. Limitations ────────────────────────────────────────────────────────
\section{Limitations}
\label{sec:limitations}

Several limitations should qualify interpretation of these findings. The keyword-based taxonomy
captures thematic \textit{presence}, not depth of treatment; a paper assigned to Category~E may
engage with ethics as a central concern or as a passing acknowledgement. Category~G's near-universal
rate partly reflects its design as a broad methodological marker. Future work should apply LLM-based
per-paper primary-category classification to produce a more discriminative taxonomy. The
international co-authorship figure~(9.9\%) now reflects OpenAlex-enriched affiliations for 96.9\%
of IEEE/ACM/ERIC records, but the 3.1\% lacking enriched affiliations may introduce minor undercount.
The 24 papers for which PDFs were not retrievable introduce a small content analysis exclusion bias
in Pipeline~2. Finally, the rapid evolution of the field means that any corpus assembled to mid-2026
will be immediately superseded; the open reproducible pipeline is designed for incremental update,
and the GitHub repository~\cite{sok_repo} provides infrastructure for future corpus extension.

% ── IX. Conclusion ────────────────────────────────────────────────────────────
\section{Conclusion}
\label{sec:conclusion}

This paper presents a Systematization of Knowledge review of the GenAI-in-assessment literature
within engineering, STEM, and computing education, drawing on 641~papers (617 full-text coded)
with bibliometric analysis spanning 2,122 deduplicated records from five databases. The field has
grown at 296.6\% per annum---approximately~$9\times$ from 2023 to 2025---dominated by conference
publications and geographically concentrated in the United States and China. Keyword co-occurrence
analysis identifies four thematic clusters: Disciplinary Contexts, GenAI/Assessment Discourse, LLM
Applications in Education, and CS/Computing Pedagogy.

The seven-category taxonomy reveals that diagnostic and normative work accounts for the preponderance
of publications, while implementation and controlled evaluation of assessment interventions remains
comparatively sparse. The central SoK contribution is the evidence quality mapping: Formative and
Adaptive Assessment~(D) has the most developed empirical base and is ready for evidence-based
adoption; Automated Grading~(A) requires validity and fairness work before high-stakes deployment;
Ethics/Policy~(E) and Authentic Redesign~(B) lack sufficient empirical foundations to support
confident practice recommendations. Fourteen priority gaps are identified, with automated grading
validity, AI-integrated assessment evaluation, detection fairness for non-native speakers, and
cross-national equity studies as the highest-urgency actionable targets. The reproducible dual-pipeline
infrastructure and open corpus~\cite{sok_repo} provide a reusable foundation for updating this
synthesis as the literature continues to evolve.

% ── References ────────────────────────────────────────────────────────────────
\bibliographystyle{IEEEtran}
\bibliography{tale2026_sok_paper}

\end{document}
